\documentclass{article}
\usepackage[margin=1.5cm]{geometry}
\usepackage{amsmath}
\usepackage{enumitem}

\setlist{nosep,leftmargin=*}

\begin{document}
\twocolumn

\title{Chapter 6 In-Class Exercise: Defining and Calling Functions}
\author{Prof. Jordan C. Hanson}
\date{}

\maketitle

\section{Defining a Function}

\begin{enumerate}

\item A function is a named block of code that performs a task.
A simple Python function is defined using \texttt{def}:

\begin{verbatim}
def square(x):
    y = x * x
    return y
\end{verbatim}

Here, \texttt{x} is a \textbf{parameter}.  The
\texttt{return} statement sends the calculated value back
to the program.

Call the function using

\begin{verbatim}
print(square(3))
print(square(5))
\end{verbatim}

Before running the code, predict both results.

\end{enumerate}


\section{Arguments and Return Values}

\begin{enumerate}

\item A value supplied when a function is called is an
\textbf{argument}.  Try

\begin{verbatim}
a = square(4)
b = square(6)

print(a)
print(b)
print(a + b)
\end{verbatim}

Answer the following:

\begin{itemize}
\item What argument is supplied in \texttt{square(4)}?
\item What value is returned?
\item What values are stored in \texttt{a} and \texttt{b}?
\end{itemize}

\end{enumerate}


\section{A Function with a Pole}

\begin{enumerate}

\item Consider the mathematical function

\[
f(x)=\frac{1}{x-1}.
\]

This function has a \textbf{pole} at $x=1$, where the
denominator becomes zero.

Define the function in Python:

\begin{verbatim}
def f(x):
    return 1 / (x - 1)
\end{verbatim}

Evaluate it at several points between 0 and 2:

\begin{verbatim}
print(f(0.0))
print(f(0.5))
print(f(0.9))
print(f(1.1))
print(f(1.5))
print(f(2.0))
\end{verbatim}

\begin{itemize}

\item What happens to $f(x)$ as $x$ approaches 1 from
the left?

\item What happens as $x$ approaches 1 from the right?

\item Predict what will happen if you execute

\begin{verbatim}
print(f(1.0))
\end{verbatim}

Why?

\end{itemize}

\end{enumerate}


\section{In-Class Exercise}

\begin{enumerate}

\item Define a new function

\[
g(x)=\frac{x+2}{x-1}.
\]

Your function should have the form

\begin{verbatim}
def g(x):
    # calculate the function
    # return the result
\end{verbatim}

\begin{itemize}

\item Evaluate \texttt{g(x)} at

\[
x=0,\;0.5,\;0.9,\;1.1,\;1.5,\;2.
\]

Use a simple \texttt{for} loop:

\begin{verbatim}
values = [0, 0.5, 0.9,
          1.1, 1.5, 2]

for x in values:
    print(x, g(x))
\end{verbatim}

\item Which statement \textbf{defines} the function?

\item Which expression \textbf{calls} the function?

\item What value is the \textbf{argument} during the first
pass through the loop?

\item What is returned by \texttt{g(2)}?

\item Where is the pole of $g(x)$?  Explain how you can
identify it directly from the denominator.

\end{itemize}

\end{enumerate}

\end{document}